% Autorska priprava k Realisticky.cz 020206; ulohy jsou autorske a obmenene.
% Zamerne nenahrazuje vyklad presnymi kopiemi cizich textu/prikladu.
\documentclass[11pt,a4paper]{article}
\input{../../sablony/pripravy/priprava.sty}
\usepackage{tikz}
\nastavPripravu{1. ročník}{Rovnice a nerovnice}{77}{Lineární nerovnice se závorkami a zlomky}
\begin{document}
\blok{Návaznost a motivace}
Rozšíříme jednoduché nerovnice o složitější úpravy. Všechny jmenovatele v základních úlohách jsou kladné konstanty, proto se po jejich odstranění směr nerovnosti nemění.
\blok{Princip a odvození}
Nejprve odstraň závorky a \emph{konstantní} jmenovatele. Při přechodu k poslednímu kroku dej zvláštní pozor na záporný koeficient u $x$. Ověř jednu hodnotu, která do řešení patří, a jednu, která nepatří.
\blok{Řešené příklady: od jednoduchého ke složitějšímu}
\begin{prikladReseniBox}[title={Příklad 1}]
Vyřešte $\frac{2x-1}{3}>\frac{x+2}{2}$.
\tcblower
\begin{align*}
2(2x-1)&>3(x+2) \\
 4x-2&>3x+6 \\
 x&>8
\end{align*}

$K=(8;+\infty)$. Násobili jsme kladným číslem $6$.
\end{prikladReseniBox}
\begin{prikladReseniBox}[title={Příklad 2}]
Vyřešte $3-2(x-1)\le\frac{x+4}{2}$.
\tcblower
\begin{align*}
6-4(x-1)&\le x+4 \\
 10-4x&\le x+4 \\
 6&\le5x \\
 x&\ge\frac65
\end{align*}

Zápis $K=\langle\frac65;+\infty)$; rovnost v krajním bodě je povolena.
\end{prikladReseniBox}
\poznamka{Věnuj čas výběru správných závorek u koncových bodů; odkazuj na práci s intervaly z 1. ročníku.}
\blok{Samostatné procvičování}
\begin{zadaniReseniBox}[title={Úloha 1}]
$2(x-1)\le3x+4$.
\tcblower
\begin{align*}
2x-2&\le3x+4 \\
 -6&\le x \\
 x&\ge-6
\end{align*}
\end{zadaniReseniBox}
\begin{zadaniReseniBox}[title={Úloha 2}]
$\frac{x+3}{4}>\frac{x-1}{2}$.
\tcblower
\begin{align*}
x+3&>2x-2 \\
 x&<5 \\
 K&=(-\infty;5)
\end{align*}
\end{zadaniReseniBox}
\begin{zadaniReseniBox}[title={Úloha 3}]
$5x-7<5x+1$.
\tcblower
\begin{align*}
-7&<1 \\
 K&=\mathbb R
\end{align*}
\end{zadaniReseniBox}
\blok{Volitelná rezerva}
\begin{zadaniReseniBox}[title={Úloha 1}]
Rozhodněte, zda řešení $-3x+2>8$ leží vlevo, nebo vpravo od $-2$.
\tcblower
\begin{align*}
-3x&>6 \\
 x&<-2 \\
 K&=(-\infty;-2)
\end{align*}

První kontrolní číslo $-3$ vyhovuje, číslo $0$ nikoliv.
\end{zadaniReseniBox}
\blok{Domácí úkol}
\begin{zadaniReseniBox}[title={Úloha 1}]
$\frac{3x+1}{2}\ge x+2$.
\tcblower
\begin{align*}
3x+1&\ge2x+4 \\
 x&\ge3 \\
 K&=\langle3;+\infty)
\end{align*}
\end{zadaniReseniBox}
\begin{zadaniReseniBox}[title={Úloha 2}]
$2-\frac{x+1}{3}<1$.
\tcblower
\begin{align*}
6-(x+1)&<3 \\
 5-x&<3 \\
 x&>2
\end{align*}
\end{zadaniReseniBox}
\begin{zadaniReseniBox}[title={Úloha 3}]
$4(x-2)\ge4x-7$.
\tcblower
\begin{align*}
4x-8&\ge4x-7 \\
 -8&\ge-7 \\
 K&=\varnothing
\end{align*}
\end{zadaniReseniBox}
\poznamka{Orientační časové rozvržení: 5 + 12 + 12 + 10 + 6 = 45 min. Přednost má pochopení jádra; rezerva není součástí časového součtu. Domácí úlohy jsou určeny na samostatnou práci, řešení zůstávají učiteli.}
\end{document}
